\section{讨论与局限}
\label{sec:discussion}

\noindent\textbf{局限与未来工作。}该层微分的是一个局部鲁棒不动点：外层GNC不是全局求解器，不
提供崩溃点保证，相干多数外点还可能使鲁棒全局最优本身出错。外部基线（Theseus-GM与
SDPRLayers+SCS）仅在合成3D--3D上评测，真实数据前端为手写的（SIFT+StereoSGBM与FPFH），夜间
In-The-Dark证据为10对窗口，且未使用学习描述子。将评测扩展到学习前端、真实数据外部基线与整夜
序列是自然的下一步；通用40根QPEP内层求解器与GPU批处理留作未来工作。

\noindent\textbf{模型崩溃。}即使$F_{\mathrm{rob}}$有完美的全局求解器，当相干外点多于内点时它
也可能返回错误位姿。一个一维例子：一个内点在$0$、两个相干外点在$M$，当$M\to\infty$时
$F(0)\to2c^2$、$F(M)\to c^2$，故全局极小是外点位姿。这一可辨识性极限需要一致性检查、内点计数
诊断或外点比例先验。

\noindent\textbf{边界残差。}$\lambda_\parallel$在$e=c^2$处最负，而不是在$e\to\infty$处。粗大
外点在$O(\|\mathbf{r}\|^{-4})$意义下被抑制，但刚好越过重降阈值的点具有最大的梯度影响，故边界
残差既是训练中最敏感的部分，也是曲率修正最关键之处。

\noindent\textbf{相对RANSAC的运行时间。}RANSAC是强且调优良好的基线，该层去掉其不可微外壳的
代价是局部鲁棒估计器。相对固定预算RANSAC，加速为$22$--$32\times$；相对自适应预算RANSAC，在
低外点比例更慢，$40\%$附近持平，$70\%$时约快$8\times$；LO-RANSAC精度持平，但比自适应RANSAC
慢$1.03$--$1.52\times$，因此比较必须注明迭代预算。

\noindent\textbf{核尺度与残差模型。}框架覆盖GM、Cauchy、Welsch与TLS，但尺度$c$至关重要：
固定$c$可能在变化噪声下失效，学习的$c$需要$\det W_i=1$之类归一化。这些推导对关于$\mathbf{t}$
仿射、关于提升四元数二次的残差是精确的；重投影误差不是多项式，QPEP在该情形使用代数代理。

\noindent\textbf{适用范围。}当匹配含粗大外点、需要穿过位姿的端到端梯度、外点比例低于崩溃区，
且$6\times6$反向与固定规模代数前向符合预算时，该层适用；当目标是全局最优性时，可认证全局求解器
仍是更合适的选择。
